\metatitle{Graphs and colored matroids}
\metadescription{Graphic matroids in two colors, prescribed-indegree
orientations, directed-cycle exchange, and circuits described by cuts.}

\section[graphsAndColoredMatroids]{Graphs and colored matroids}

Graphs enter the theory in two related ways. Spanning trees give the usual
graphic matroid in two colors. With more colors, fixed-indegree orientations
give a family in which symmetric cycle exchange is literally reversal of a
directed cycle.

\subsection[graphicMatroids]{The classical graphic matroid}

Let $G=(V,E)$ be a finite connected graph. Its graphic matroid $M(G)$ has
edge set $E$; its bases are the spanning trees, its independent sets are the
forests, and its circuits are the edge sets of graph cycles. The circuits of
the dual matroid are the bonds, that is, the minimal nonempty edge cuts.

For every spanning tree $T$, color its edges $1$ and the remaining edges
$0$. This gives the binary word
\[
 u_T(e)=
 \begin{cases}
 1,&e\in T,\\
 0,&e\notin T.
 \end{cases}
\]

\begin{proposition}[The graphic matroid in two colors]
The words $u_T$, as $T$ ranges over the spanning trees of $G$, form a
two-color colored matroid. Its color-$1$ projection is $M(G)$, and its
color-$0$ projection is the cographic matroid $M(G)^*$.

If a partial assignment colors the edges of $I$ by $1$ and the edges of $J$
by $0$, then it is independent if and only if $I$ is a forest and
$G\setminus J$ is connected. Its monochromatic circuits are graph cycles in
color $1$ and bonds in color $0$.
\end{proposition}

\begin{proof}
The spanning trees are the bases of $M(G)$, so the two-color result gives
symmetric cycle exchange and identifies the two projections as $M(G)$ and
$M(G)^*$.

A set $I$ is independent in $M(G)$ precisely when it is a forest. A set $J$
is independent in $M(G)^*$ precisely when it is contained in the complement
of a spanning tree, which is equivalent to $G\setminus J$ containing a
spanning tree. The latter condition says that $G\setminus J$ is connected.
The description of the circuits follows from the ordinary graphic and
cographic matroids.
\end{proof}

Thus the usual graphic matroid is fully present: forests are independent,
spanning trees are bases, graph cycles are circuits, and bonds are
cocircuits. The second color records the dual at the same time.

\subsection[fixedIndegreeOrientations]{Fixed-indegree orientations}

Now let the vertex set of a finite loopless multigraph $G$ be the set of
colors. Each edge is a position and may receive either of its endpoint
colors. Equivalently, we orient the edge toward the vertex giving its color.
Fix nonnegative integers $r_v$, for $v\in V$, with
\[
 \sum_{v\in V}r_v=|E|,
\]
and take all orientations with indegree $r_v$ at every vertex $v$.

\begin{proposition}[Directed-cycle exchange]
Every nonempty fixed-indegree orientation family is a colored matroid. For
two orientations and any edge on which they disagree, symmetric cycle
exchange is obtained by reversing a directed cycle containing that edge.
\end{proposition}

\begin{proof}
Keep the edges on which the two orientations disagree and orient them as in
the first orientation. At every vertex, the indegree equals the outdegree:
the two original orientations have the same indegree there. Every edge of a
finite balanced directed graph lies on a directed cycle. Reverse such a
cycle through the prescribed edge in both orientations. All indegrees remain
unchanged, so both resulting orientations belong to the family.
\end{proof}

\begin{example}[The triangle]
Order the edges of a triangle as
\[
 a=12,\qquad b=23,\qquad c=31.
\]
There are two orientations with indegree one at every vertex. Their words in
the order $(a,b,c)$ are
\[
 \mathtt{231}\qquad\text{and}\qquad\mathtt{123}.
\]
They are the two directed cycles, and symmetric cycle exchange reverses all
three edges.
\end{example}

\subsection[orientationIndependence]{Independence and cuts}

Fix a nonempty orientation family. A partial orientation is independent if
it can be completed to the prescribed indegrees. This has a cut description.

For $X\subseteq V$, let $E(X)$ be the set of edges with both endpoints in
$X$, and let $\delta(X)$ be the set of edges with one endpoint in $X$. If
$\tau$ is a partial orientation, let $i_\tau(X)$ be the number of already
oriented edges of $\delta(X)$ that point into $X$.

\begin{proposition}[Cut criterion]
A partial orientation $\tau$ extends to an orientation with indegrees
$(r_v)_{v\in V}$ if and only if
\[
 i_\tau(X)
 \leq
 \sum_{v\in X}r_v-|E(X)|
 \qquad\text{for every }X\subseteq V.
\]
\end{proposition}

\begin{proof}
Let $F$ be the set of edges already oriented and let $R=E\setminus F$. At a
vertex $v$, the remaining edges must contribute
\[
 b_v=r_v-d^-_\tau(v)
\]
incoming edges. The cut inequality for $X=\{v\}$ gives $b_v\geq0$, and
$\sum_{v\in V}b_v=|R|$. Assigning every edge of $R$ to its chosen head is a matching
problem: make one vertex for each edge of $R$, and make $b_v$ slots at every
vertex $v$; an edge may be matched to a slot at either endpoint.

Hall's condition for this matching is
\[
 \sum_{v\in X}b_v\geq |R(X)|
 \qquad\text{for every }X\subseteq V.
\]
Indeed, every collection of remaining edges lies inside the subgraph on the
union of its endpoints. Now
\[
 \sum_{v\in X}d^-_\tau(v)
 =|F\cap E(X)|+i_\tau(X),
 \qquad
 |R(X)|=|E(X)|-|F\cap E(X)|.
\]
Substitution gives exactly the displayed cut inequalities.
\end{proof}

Assigning an edge a color other than its endpoints gives a one-element
circuit. The remaining circuits are partial orientations described by cuts.

\begin{corollary}[Circuits are minimal overloaded cuts]
Let $C$ be a circuit consisting of endpoint assignments. There is a set
$X\subseteq V$ such that every edge prescribed by $C$ belongs to $\delta(X)$
and points into $X$, and
\[
 |C|=\sum_{v\in X}r_v-|E(X)|+1.
\]
Conversely, every inclusion-minimal partial orientation violating a cut
inequality is a circuit.
\end{corollary}

\begin{proof}
Since $C$ does not extend, it violates the cut inequality for some $X$. If
one of its assignments were not an edge directed across the cut into $X$,
removing that assignment would leave the same violation, contrary to
minimality. Removing any assignment must repair the inequality, so the left
side exceeds the right side by exactly one. The converse follows from the
cut criterion.
\end{proof}

For the triangle with indegrees $(1,1,1)$, directing both edges incident to
one vertex into that vertex is a circuit: its prescribed indegree is already
two. Directing both edges away from one vertex is another circuit, detected
by the cut consisting of the other two vertices. The bases are the two
directed cycles, the exchanges are directed cycles, and the circuits are
minimal cut overloads.

\subsection[treePartitions]{Why spanning-tree partitions are different}

It is tempting to require every color class to be a spanning tree. This does
not always satisfy symmetric cycle exchange, even with two colors.

\begin{example}[Two spanning trees]
In $K_4$, let
\[
 T=\{12,23,34\},
 \qquad
 T'=\{12,13,24\}.
\]
Both $T$ and $T'$ have spanning-tree complements. Consider the two binary
words whose color-$1$ classes are $T$ and $T'$, and prescribe exchange at
$34\in T\setminus T'$. The possible incoming edges from $T'$ are $13$ and
$24$.

Replacing $34$ by $13$ produces the triangle $\{12,13,23\}$. Replacing it
by $24$ produces the tree $\{12,23,24\}$, but its complement
$\{13,14,34\}$ is a triangle. Neither exchange gives two valid ordered
partitions into spanning trees.
\end{example}

The fixed-indegree orientation construction succeeds because its feasibility
conditions are preserved by reversing directed cycles. Constructions that
encode an arbitrary ordinary graphic matroid using more colors are described
on the \href{constructions.htm}{constructions and examples} page.
